% LuaLaTeX でコンパイルしてください
\documentclass[a4paper,11pt,ja=standard]{bxjsarticle}

% パッケージの読み込み
\usepackage{luatexja}
\usepackage{luatexja-fontspec}
\usepackage{amsmath, amssymb, amsthm}
\usepackage{tikz}
\usetikzlibrary{arrows.meta, positioning, calc, decorations.pathreplacing, shapes.geometric, cd}
\usepackage{tcolorbox}
\tcbuselibrary{theorems, skins, breakable}
\usepackage{hyperref}
\usepackage{listings}
\usepackage{xcolor}
\usepackage{enumitem}

% 日本語フォントの設定（適宜変更してください）

% Leanコードのスタイル設定
\lstdefinestyle{leanstyle}{
    language=ML,
    basicstyle=\small\ttfamily,
    keywordstyle=\color{blue}\bfseries,
    commentstyle=\color{green!60!black},
    stringstyle=\color{red},
    showstringspaces=false,
    breaklines=true,
    frame=single,
    numbers=left,
    numberstyle=\tiny\color{gray},
    backgroundcolor=\color{gray!10}
}

% 定理環境の設定
\theoremstyle{definition}
\newtheorem{definition}{定義}[section]
\newtheorem{theorem}[definition]{定理}
\newtheorem{proposition}[definition]{命題}
\newtheorem{lemma}[definition]{補題}
\newtheorem{example}[definition]{例}
\newtheorem{remark}[definition]{注意}
\newtheorem{corollary}[definition]{系}
\newtheorem{exercise}[definition]{演習問題}

% ボックス環境の設定
\newtcolorbox{keyidea}{
    enhanced,
    breakable,
    colback=blue!5!white,
    colframe=blue!75!black,
    title=重要なアイディア,
    fonttitle=\bfseries\sffamily,
    attach boxed title to top left={yshift=-2mm, xshift=5mm},
    boxed title style={colback=blue!75!black}
}

\newtcolorbox{intuition}{
    enhanced,
    breakable,
    colback=purple!5!white,
    colframe=purple!70!black,
    title=直感的理解,
    fonttitle=\bfseries\sffamily,
    attach boxed title to top left={yshift=-2mm, xshift=5mm},
    boxed title style={colback=purple!70!black}
}

\newtcolorbox{geometric}{
    enhanced,
    breakable,
    colback=orange!5!white,
    colframe=orange!80!black,
    title=幾何学的直感,
    fonttitle=\bfseries\sffamily,
    attach boxed title to top left={yshift=-2mm, xshift=5mm},
    boxed title style={colback=orange!80!black}
}

\newtcolorbox{bourbaki}{
    enhanced,
    breakable,
    colback=green!5!white,
    colframe=green!60!black,
    title=Bourbakiの視点,
    fonttitle=\bfseries\sffamily,
    attach boxed title to top left={yshift=-2mm, xshift=5mm},
    boxed title style={colback=green!60!black}
}

% タイトル情報
\title{構造塔の理論と具体例\\{\large Bourbaki的母構造への応用}}
\author{%
    Lean 形式化: CODEX (OpenAI)\\
    \LaTeX 解説: Claude Code (Anthropic)
}
\date{\today}

\begin{document}

\maketitle

\begin{abstract}
本稿では、Nicolas Bourbakiが提唱した「母構造」(structures mères) の概念を現代的な構造塔 (Structure Tower) の理論によって具体化し、位相構造、代数構造、順序構造、およびそれらの組み合わせにおける豊富な具体例を通じて解説する。各構造に対して、厳密な数学的定義とLean 4による形式化、さらに図式による視覚的理解を提供する。本解説は学部2年生以上を対象とし、抽象的な数学理論と具体的な数学対象を橋渡しすることを目指す。
\end{abstract}

\tableofcontents
\newpage

\section{導入：なぜ構造塔を学ぶのか}

\subsection{Bourbakiの夢：数学の統一}

20世紀、Nicolas Bourbakiというペンネームを持つ数学者集団は、数学全体を「構造」という概念で統一しようとした。彼らは、すべての数学的対象は以下の三つの「母構造」のいずれか、または組み合わせとして理解できると考えた：

\begin{enumerate}
    \item \textbf{代数構造} (algebraic structures)：群、環、体など
    \item \textbf{順序構造} (order structures)：順序集合、束など
    \item \textbf{位相構造} (topological structures)：位相空間、距離空間など
\end{enumerate}

\begin{center}
\begin{tikzpicture}[scale=1.2]
    % 3つの円
    \draw[thick, fill=red!20, opacity=0.6] (0,0) circle (2cm) node[above, yshift=1.5cm] {\textbf{代数構造}};
    \draw[thick, fill=blue!20, opacity=0.6] (2.5,0) circle (2cm) node[above, yshift=1.5cm] {\textbf{位相構造}};
    \draw[thick, fill=green!20, opacity=0.6] (1.25,-2) circle (2cm) node[below, yshift=-1.5cm] {\textbf{順序構造}};

    % 交わり
    \node at (1.25,0.3) {位相群};
    \node at (0.5,-1) {順序群};
    \node at (2,-1) {順序位相};
    \node at (1.25,-0.8) {\scriptsize 位相順序群};

    % 個別の例
    \node[red!80!black] at (-1.2,0.5) {環};
    \node[red!80!black] at (-0.8,-0.5) {体};
    \node[blue!80!black] at (3.7,0.5) {多様体};
    \node[blue!80!black] at (3.3,-0.5) {距離空間};
    \node[green!80!black] at (1.25,-3) {束};
    \node[green!80!black] at (0.3,-2.5) {フィルター};
\end{tikzpicture}
\end{center}

\begin{bourbaki}
構造塔の理論は、Bourbakiの構造主義を形式化したものである。異なる数学分野で現れる「階層的包含関係」を統一的に扱うことができる。
\end{bourbaki}

\subsection{構造塔とは何か}

構造塔は、数学的対象が持つ「段階的な豊かさ」を形式化する。

\begin{center}
\begin{tikzpicture}[scale=1.3]
    % 塔の描画
    \foreach \i/\name/\color/\desc in {
        0/集合/blue!15/{要素と属性},
        1/群/green!20/{二項演算+結合律+単位元+逆元},
        2/環/orange!25/{加法群+乗法+分配律},
        3/体/red!30/{環+乗法逆元}
    } {
        \fill[\color] (-3,\i*1.2) rectangle (3,\i*1.2+1);
        \draw[thick] (-3,\i*1.2) rectangle (3,\i*1.2+1);
        \node[font=\bfseries] at (0,\i*1.2+0.7) {\name};
        \node[font=\scriptsize] at (0,\i*1.2+0.3) {\desc};
    }

    % 矢印と説明
    \draw[-Stealth, ultra thick, blue] (4,4) -- (4,0.5);
    \node[right, align=left, blue, font=\small] at (4.3,2.3) {
        \textbf{より豊かな構造}\\
        より多くの性質\\
        より強い制約
    };

    % 包含関係
    \draw[-Stealth, thick, red, dashed] (1.5,0.5) to[out=45, in=-45] (1.5,1.7);
    \draw[-Stealth, thick, red, dashed] (1.5,1.7) to[out=45, in=-45] (1.5,2.9);
    \draw[-Stealth, thick, red, dashed] (1.5,2.9) to[out=45, in=-45] (1.5,4.1);
    \node[right, red, font=\small] at (2,2.5) {単調性};
\end{tikzpicture}
\end{center}

\begin{keyidea}
構造塔の本質は\textbf{単調性}にある：
\[
\text{下の層に属する} \implies \text{上の層にも属する}
\]
つまり、「一度満たした性質は失われない」という数学的直感を形式化したものである。
\end{keyidea}

\section{構造塔の基礎理論}

\subsection{形式的定義}

\begin{definition}[構造塔 (Structure Tower with Minimum Layer)]
構造塔 $T$ は以下のデータと公理からなる：

\textbf{データ：}
\begin{enumerate}
    \item \textbf{台集合} $X = T.\text{carrier}$ : すべての要素を含む型
    \item \textbf{添字集合} $I = T.\text{Index}$ : 前順序 $(I, \leq)$ が入った型
    \item \textbf{層関数} $\lambda : I \to \mathcal{P}(X)$, $i \mapsto T.\text{layer}(i)$
    \item \textbf{最小層関数} $\mu : X \to I$, $x \mapsto T.\text{minLayer}(x)$
\end{enumerate}

\textbf{公理：}
\begin{enumerate}
    \item \textbf{被覆性} (Covering)：
    \[
    \forall x \in X, \exists i \in I, \quad x \in \lambda(i)
    \]
    すべての要素はどこかの層に属する。

    \item \textbf{単調性} (Monotonicity)：
    \[
    \forall i, j \in I, \quad i \leq j \implies \lambda(i) \subseteq \lambda(j)
    \]
    添字が大きくなれば層も大きくなる。

    \item \textbf{最小層の帰属} (Minimum Layer Membership)：
    \[
    \forall x \in X, \quad x \in \lambda(\mu(x))
    \]
    各要素はその最小層に属する。

    \item \textbf{最小性} (Minimality)：
    \[
    \forall x \in X, \forall i \in I, \quad x \in \lambda(i) \implies \mu(x) \leq i
    \]
    最小層は本当に最小である。
\end{enumerate}
\end{definition}

\begin{intuition}
\begin{itemize}
    \item \textbf{被覆性}：「どの要素も宙に浮いていない」
    \item \textbf{単調性}：「性質は累積的に増える」
    \item \textbf{最小層}：「その要素が初めて現れる階層」
    \item \textbf{最小性}：「それより下の層には属さない」
\end{itemize}
\end{intuition}

\subsection{構造塔の図式表現}

構造塔の要素 $x$ がどのように層を上昇するかを図示する：

\begin{center}
\begin{tikzpicture}[scale=1.4]
    % 層の描画
    \foreach \i in {0,1,2,3,4} {
        \draw[thick, fill=blue!10] (-2.5,\i*0.8) rectangle (2.5,\i*0.8+0.7);
        \node[left, font=\small] at (-2.7,\i*0.8+0.35) {$i_{\i}$};
    }

    % 要素xの軌跡
    \node[circle, fill=red, inner sep=3pt, label=right:{$x \notin \lambda(i_0)$}] at (0,0.35) {};
    \node[circle, fill=red, inner sep=3pt, label=right:{$x \notin \lambda(i_1)$}] at (0,1.15) {};
    \node[circle, fill=green!70!black, inner sep=3pt, label=right:{$x \in \lambda(i_2) \leftarrow \mu(x) = i_2$}] (x2) at (0,1.95) {};
    \node[circle, fill=green!70!black, inner sep=3pt, label=right:{$x \in \lambda(i_3)$}] at (0,2.75) {};
    \node[circle, fill=green!70!black, inner sep=3pt, label=right:{$x \in \lambda(i_4)$}] at (0,3.55) {};

    % 注釈
    \draw[decorate, decoration={brace, amplitude=10pt, mirror}, thick, red]
        (-2.5,-0.1) -- (-2.5,1.5) node[midway, left, xshift=-10pt, align=center, font=\small] {$x$は\\属さない};
    \draw[decorate, decoration={brace, amplitude=10pt, mirror}, thick, green!70!black]
        (-2.5,1.7) -- (-2.5,4.0) node[midway, left, xshift=-10pt, align=center, font=\small] {$x$は\\属する};

    % 最小層の強調
    \draw[ultra thick, orange, dashed] (-2.7,1.6) -- (2.7,1.6);
    \node[orange, right, font=\small\bfseries] at (2.8,1.95) {最小層 $\mu(x)$};
\end{tikzpicture}
\end{center}

\section{具体例 I：位相構造}

\subsection{離散位相の構造塔}

最も単純な構造塔の例から始める。

\begin{example}[離散位相塔]
位相空間 $X$ に対して、以下のように構造塔を定義する：
\begin{align*}
T.\text{carrier} &:= X \\
T.\text{Index} &:= \mathcal{P}(X) \quad \text{（部分集合の族）} \\
T.\text{layer}(F) &:= F \quad \text{（集合そのもの）} \\
T.\text{minLayer}(x) &:= \{x\} \quad \text{（単集合）}
\end{align*}
ここで、$\mathcal{P}(X)$ の順序は包含関係 $\subseteq$ とする。
\end{example}

\begin{geometric}
離散位相塔では、各層は単に部分集合であり、最小層は「その点だけを含む単集合」である。これは最も原始的な構造塔である。
\end{geometric}

\subsubsection{公理の検証}

離散位相塔が構造塔の公理を満たすことを確認しよう。

\begin{proof}[公理の検証]
\begin{enumerate}
    \item \textbf{被覆性}：任意の $x \in X$ に対して、$x \in \{x\} = \lambda(\{x\})$ なので、$i = \{x\}$ を取れば $x \in \lambda(i)$。

    \item \textbf{単調性}：$F \subseteq G$ とする。すると $\lambda(F) = F \subseteq G = \lambda(G)$。

    \item \textbf{最小層の帰属}：$\mu(x) = \{x\}$ なので、$x \in \{x\} = \lambda(\mu(x))$。

    \item \textbf{最小性}：$x \in \lambda(F) = F$ とする。すると $x \in F$ より $\{x\} \subseteq F$、つまり $\mu(x) = \{x\} \leq F$。
\end{enumerate}
\end{proof}

\subsection{位相の細かさによる構造塔}

より高度な例として、位相の包含関係による構造塔を考える。

\begin{example}[位相の細かさ塔]
集合 $X$ 上の位相全体を $\text{Top}(X)$ とする。位相 $\tau_1 \leq \tau_2$ を「$\tau_1$ は $\tau_2$ より粗い」と定義する（つまり $\tau_1 \subseteq \tau_2$ as sets of open sets）。

構造塔を以下のように定義：
\begin{align*}
T.\text{carrier} &:= \mathcal{P}(X) \quad \text{（$X$の部分集合全体）} \\
T.\text{Index} &:= \text{OrderDual}(\text{Top}(X)) \\
T.\text{layer}(\tau) &:= \{U \subseteq X \mid U \text{ は } \tau\text{-開集合}\} \\
T.\text{minLayer}(U) &:= \text{生成位相}(\{U\})
\end{align*}
\end{example}

\begin{intuition}
この構造塔では、層は「ある位相で開集合となる部分集合の族」である。位相が細かくなるほど（つまり開集合が増えるほど）、層が大きくなる。
\end{intuition}

\begin{center}
\begin{tikzpicture}[scale=1.3]
    % 位相の階層
    \node[draw, thick, fill=blue!10, minimum width=8cm, minimum height=0.8cm] (top) at (0,4)
        {離散位相 $\tau_{\text{discrete}}$：すべての部分集合が開};
    \node[draw, thick, fill=blue!20, minimum width=7cm, minimum height=0.8cm] (mid2) at (0,3)
        {有限補位相 $\tau_{\text{cofinite}}$};
    \node[draw, thick, fill=blue!30, minimum width=6cm, minimum height=0.8cm] (mid1) at (0,2)
        {生成位相 $\tau_{\{U\}}$};
    \node[draw, thick, fill=blue!40, minimum width=5cm, minimum height=0.8cm] (bottom) at (0,1)
        {密着位相 $\tau_{\text{indiscrete}} = \{\emptyset, X\}$};

    % 矢印
    \draw[-Stealth, thick, green!70!black] (bottom) -- (mid1) node[midway, right, font=\scriptsize] {細かくなる};
    \draw[-Stealth, thick, green!70!black] (mid1) -- (mid2);
    \draw[-Stealth, thick, green!70!black] (mid2) -- (top);

    % 注釈
    \node[left, font=\small, align=right] at (-4.5,1) {最も粗い\\開集合が少ない};
    \node[left, font=\small, align=right] at (-4.5,4) {最も細かい\\開集合が多い};
\end{tikzpicture}
\end{center}

\section{具体例 II：代数構造}

\subsection{主イデアル環の構造塔}

代数学の基本的対象である環とそのイデアルを構造塔で捉える。

\begin{example}[主イデアル塔]
可換環 $R$ に対して、以下のように構造塔を定義する：
\begin{align*}
T.\text{carrier} &:= R \quad \text{（環の元）} \\
T.\text{Index} &:= \text{Ideal}(R) \quad \text{（イデアル全体）} \\
T.\text{layer}(I) &:= I \quad \text{（イデアルを集合とみなす）} \\
T.\text{minLayer}(x) &:= \langle x \rangle = Rx \quad \text{（主イデアル）}
\end{align*}
ここで、イデアルの順序は包含関係 $\subseteq$ とする。
\end{example}

\subsubsection{整数環での具体例}

最も基本的な例として整数環 $\mathbb{Z}$ を考える。

\begin{center}
\begin{tikzpicture}[scale=1.4]
    % イデアルの階層
    \node[draw, thick, fill=red!10, minimum width=7cm, minimum height=0.7cm] (top) at (0,4)
        {$\langle 1 \rangle = \mathbb{Z}$ (全体)};
    \node[draw, thick, fill=red!20, minimum width=6cm, minimum height=0.7cm] (i2) at (0,3.2)
        {$\langle 2 \rangle = 2\mathbb{Z}$ (偶数)};
    \node[draw, thick, fill=red!30, minimum width=5cm, minimum height=0.7cm] (i4) at (0,2.4)
        {$\langle 4 \rangle = 4\mathbb{Z}$};
    \node[draw, thick, fill=red!40, minimum width=4cm, minimum height=0.7cm] (i8) at (0,1.6)
        {$\langle 8 \rangle = 8\mathbb{Z}$};
    \node[draw, thick, fill=red!50, minimum width=3cm, minimum height=0.7cm] (i16) at (0,0.8)
        {$\langle 16 \rangle = 16\mathbb{Z}$};
    \node[draw, thick, fill=red!60, minimum width=2cm, minimum height=0.7cm] (zero) at (0,0)
        {$\langle 0 \rangle = \{0\}$};

    % 要素の追跡
    \node[circle, fill=blue, inner sep=2pt, label=right:{$8 \in \langle 8 \rangle$}] at (2,1.6) {};
    \draw[-Stealth, thick, dashed, blue] (2,1.6) -- (2,2.35);
    \node[circle, fill=blue, inner sep=2pt, label=right:{$8 \in \langle 4 \rangle$}] at (2,2.4) {};
    \draw[-Stealth, thick, dashed, blue] (2,2.4) -- (2,3.15);
    \node[circle, fill=blue, inner sep=2pt, label=right:{$8 \in \langle 2 \rangle$}] at (2,3.2) {};
    \draw[-Stealth, thick, dashed, blue] (2,3.2) -- (2,3.95);
    \node[circle, fill=blue, inner sep=2pt, label=right:{$8 \in \mathbb{Z}$}] at (2,4) {};

    % 最小層の強調
    \draw[ultra thick, green!70!black, dashed] (-3.7,1.3) -- (3.7,1.3);
    \node[green!70!black, left, font=\small] at (-3.8,1.6) {$\mu(8) = \langle 8 \rangle$};
\end{tikzpicture}
\end{center}

\begin{example}[具体的な計算]
整数環 $\mathbb{Z}$ において：
\begin{itemize}
    \item $6 \in \langle 3 \rangle$ である。なぜなら $6 = 3 \cdot 2 \in 3\mathbb{Z}$。
    \item $6 \in \langle 2 \rangle$ である。なぜなら $6 = 2 \cdot 3 \in 2\mathbb{Z}$。
    \item $6 \in \langle 1 \rangle = \mathbb{Z}$ である（自明）。
    \item しかし $\mu(6) = \langle 6 \rangle$ であり、$\langle 6 \rangle \subsetneq \langle 3 \rangle, \langle 2 \rangle$。
\end{itemize}
\end{example}

\subsection{部分群の構造塔}

群論における基本概念を構造塔で捉える。

\begin{example}[部分群塔]
群 $G$ に対して、以下のように構造塔を定義する：
\begin{align*}
T.\text{carrier} &:= G \quad \text{（群の元）} \\
T.\text{Index} &:= \text{Subgroup}(G) \quad \text{（部分群全体）} \\
T.\text{layer}(H) &:= H \quad \text{（部分群を集合とみなす）} \\
T.\text{minLayer}(g) &:= \langle g \rangle \quad \text{（$g$で生成される巡回部分群）}
\end{align*}
ここで、部分群の順序は包含関係 $\subseteq$ とする。
\end{example}

\begin{center}
\begin{tikzpicture}[scale=1.3]
    % 部分群の階層（例：Z/8Zの部分群）
    \node[draw, thick, fill=green!10, ellipse, minimum width=6cm, minimum height=0.8cm] (top) at (0,3.5)
        {$G$ (全体群)};
    \node[draw, thick, fill=green!20, ellipse, minimum width=5cm, minimum height=0.8cm] (h1) at (-1.5,2.3)
        {部分群 $H_1$};
    \node[draw, thick, fill=green!20, ellipse, minimum width=5cm, minimum height=0.8cm] (h2) at (1.5,2.3)
        {部分群 $H_2$};
    \node[draw, thick, fill=green!30, ellipse, minimum width=3.5cm, minimum height=0.8cm] (cyclic) at (0,1.1)
        {巡回部分群 $\langle g \rangle$};
    \node[draw, thick, fill=green!40, ellipse, minimum width=2cm, minimum height=0.8cm] (triv) at (0,0)
        {$\{e\}$ (単位群)};

    % 包含関係
    \draw[-Stealth, thick] (triv) -- (cyclic);
    \draw[-Stealth, thick] (cyclic) -- (h1);
    \draw[-Stealth, thick] (cyclic) -- (h2);
    \draw[-Stealth, thick] (h1) -- (top);
    \draw[-Stealth, thick] (h2) -- (top);

    % 要素の追跡
    \node[circle, fill=red, inner sep=2pt, label=left:{$g$}] at (-0.3,1.1) {};
    \node[font=\small, align=left] at (4,1.8) {$g$ の軌跡：\\
        $g \in \langle g \rangle \subseteq H_1 \subseteq G$};
\end{tikzpicture}
\end{center}

\begin{exercise}
巡回群 $\mathbb{Z}/12\mathbb{Z}$ のすべての部分群を列挙し、構造塔として図示せよ。特に、元 $\overline{4}$ の最小層を求めよ。
\end{exercise}

\section{具体例 III：順序構造}

\subsection{下集合による構造塔}

順序集合における基本的な構造を捉える。

\begin{example}[下集合塔]
半順序集合 $(X, \leq)$ に対して、以下のように構造塔を定義する：
\begin{align*}
T.\text{carrier} &:= X \\
T.\text{Index} &:= X \quad \text{（添字も同じ集合）} \\
T.\text{layer}(x) &:= \downarrow\! x := \{y \in X \mid y \leq x\} \quad \text{（下集合）} \\
T.\text{minLayer}(x) &:= x
\end{align*}
\end{example}

\begin{center}
\begin{tikzpicture}[scale=1.5]
    % Hasse diagram
    \node[circle, draw, fill=blue!20] (a) at (0,0) {$a$};
    \node[circle, draw, fill=blue!20] (b) at (-1,1) {$b$};
    \node[circle, draw, fill=blue!20] (c) at (1,1) {$c$};
    \node[circle, draw, fill=blue!20] (d) at (0,2) {$d$};

    \draw[thick] (a) -- (b) -- (d);
    \draw[thick] (a) -- (c) -- (d);

    % 下集合の視覚化
    \begin{scope}[xshift=5cm]
        \node[font=\bfseries] at (0,2.5) {下集合の構造};

        % ↓a
        \draw[thick, fill=red!10] (-1.5,-0.3) rectangle (1.5,0.3);
        \node at (0,0) {$\downarrow\! a = \{a\}$};

        % ↓b
        \draw[thick, fill=orange!10] (-1.5,0.7) rectangle (1.5,1.3);
        \node at (0,1) {$\downarrow\! b = \{a, b\}$};

        % ↓d
        \draw[thick, fill=green!10] (-1.5,1.7) rectangle (1.5,2.3);
        \node at (0,2) {$\downarrow\! d = \{a,b,c,d\}$};

        % 矢印
        \draw[-Stealth, thick, dashed] (1.7,0.3) -- (1.7,0.7);
        \draw[-Stealth, thick, dashed] (1.7,1.3) -- (1.7,1.7);
    \end{scope}
\end{tikzpicture}
\end{center}

\subsection{フィルターの構造塔}

位相空間論や解析学で重要な概念であるフィルターも構造塔で捉えられる。

\begin{example}[フィルター塔]
集合 $X$ 上のフィルター全体を $\text{Filter}(X)$ とする。フィルターの順序は包含関係の逆：$F \leq G$ を $F \supseteq G$ と定義する（$G$ の方が細かい）。

構造塔を以下のように定義：
\begin{align*}
T.\text{carrier} &:= \mathcal{P}(X) \\
T.\text{Index} &:= \text{OrderDual}(\text{Filter}(X)) \\
T.\text{layer}(F) &:= F \quad \text{（フィルターに属する集合）} \\
T.\text{minLayer}(S) &:= \text{principal filter } \{T \subseteq X \mid S \subseteq T\}
\end{align*}
\end{example}

\begin{intuition}
フィルターは「大きな集合の族」を形式化したもの。principal filter は「ある集合 $S$ を含むすべての集合」からなる最も粗いフィルター。
\end{intuition}

\section{具体例 IV：母構造の組み合わせ}

Bourbakiの真骨頂は、複数の母構造を組み合わせることにある。

\subsection{順序位相空間}

順序構造と位相構造を同時に持つ空間。

\begin{example}[順序位相空間の構造塔]
位相空間 $(X, \tau)$ かつ半順序集合 $(X, \leq)$ に対して：
\begin{align*}
T.\text{layer}(x) &:= \overline{\downarrow\! x} \quad \text{（下集合の閉包）}
\end{align*}
これは順序と位相の両方を反映した層である。
\end{example}

\begin{center}
\begin{tikzpicture}[scale=1.3]
    % 実数直線
    \draw[-Stealth, thick] (-3,0) -- (3,0) node[right] {$\mathbb{R}$};
    \foreach \x/\label in {-2/-2, -1/-1, 0/0, 1/1, 2/2} {
        \draw (\x,0.1) -- (\x,-0.1) node[below] {$\label$};
    }

    % 点 x = 1
    \node[circle, fill=red, inner sep=2.5pt, label=above:{$x = 1$}] at (1,0) {};

    % 下集合 (-∞, 1]
    \draw[ultra thick, blue] (-3,0.5) -- (1,0.5);
    \node[blue, above, font=\small] at (-1,0.5) {$\downarrow\! 1 = (-\infty, 1]$};

    % 閉包（この場合同じ）
    \draw[ultra thick, green!70!black, dashed] (-3,-0.5) -- (1,-0.5);
    \node[green!70!black, below, font=\small] at (-1,-0.5) {$\overline{\downarrow\! 1} = (-\infty, 1]$};

    \node[font=\small, align=left] at (0,-1.5) {
        実数の通常位相では\\
        $\overline{\downarrow\! x} = (-\infty, x]$
    };
\end{tikzpicture}
\end{center}

\subsection{位相群}

群構造と位相構造を両立させた基本的対象。

\begin{example}[位相群の構造塔]
位相群 $(G, \tau, \cdot)$ に対して、単位元 $e$ の近傍系 $\mathcal{N}(e)$ から構造塔を作る：
\begin{align*}
T.\text{layer}(U) &:= \overline{\langle U \rangle} \quad \text{（$U$で生成される部分群の閉包）}
\end{align*}
ここで $U \in \mathcal{N}(e)$ は単位元の近傍。
\end{example}

\begin{center}
\begin{tikzpicture}[scale=1.3]
    % 円（位相群の視覚化、例えば S^1）
    \draw[thick] (0,0) circle (2cm);
    \node[circle, fill=red, inner sep=2.5pt, label=above:{$e$}] at (0,2) {};

    % 近傍 U
    \draw[ultra thick, blue, ->] (0,2) arc (90:60:2cm);
    \draw[ultra thick, blue, ->] (0,2) arc (90:120:2cm);
    \node[blue, above right, font=\small] at (1,2.3) {近傍 $U$};

    % 生成される部分群
    \foreach \angle in {0,30,60,...,330} {
        \node[circle, fill=green!70!black, inner sep=1.5pt] at (\angle:2cm) {};
    }
    \node[green!70!black, right, font=\small] at (2.5,0) {$\langle U \rangle$};

    \node[font=\small, align=center] at (0,-3) {
        位相群 $S^1 = \mathbb{R}/\mathbb{Z}$\\
        単位元近傍から群全体が生成される
    };
\end{tikzpicture}
\end{center}

\begin{exercise}
$\mathbb{R}^n$ にユークリッド位相と加法群構造を入れた位相群を考える。原点の開球 $B_\epsilon(0)$ で生成される部分群を求めよ。
\end{exercise}

\section{構造塔間の射}

構造塔だけでなく、構造塔の間の「構造を保つ写像」も重要である。

\begin{definition}[構造塔の射]
構造塔 $T_1, T_2$ の間の射 $f : T_1 \to T_2$ は以下のデータからなる：
\begin{enumerate}
    \item \textbf{台写像} $f_0 : T_1.\text{carrier} \to T_2.\text{carrier}$
    \item \textbf{添字写像} $f_\# : T_1.\text{Index} \to T_2.\text{Index}$
\end{enumerate}
で、以下を満たす：
\begin{enumerate}
    \item \textbf{添字単調性}：$i \leq j \implies f_\#(i) \leq f_\#(j)$
    \item \textbf{層保存性}：$x \in T_1.\text{layer}(i) \implies f_0(x) \in T_2.\text{layer}(f_\#(i))$
    \item \textbf{最小層保存}：$f_\#(T_1.\text{minLayer}(x)) = T_2.\text{minLayer}(f_0(x))$
\end{enumerate}
\end{definition}

\begin{center}
\begin{tikzpicture}[scale=1.3]
    % T1
    \begin{scope}[xshift=-3cm]
        \node[above, font=\bfseries] at (0,3.5) {$T_1$};
        \foreach \i in {0,1,2,3} {
            \draw[thick, fill=blue!15] (-1.2,\i*0.8) rectangle (1.2,\i*0.8+0.7);
        }
        \node[circle, fill=red, inner sep=2pt, label=left:{$x$}] (x1) at (0,1.4) {};
    \end{scope}

    % 射の矢印
    \draw[-Stealth, ultra thick, purple] (-1.5,1.5) -- (1.5,1.5) node[midway, above] {$f$};

    % T2
    \begin{scope}[xshift=3cm]
        \node[above, font=\bfseries] at (0,3.5) {$T_2$};
        \foreach \i in {0,1,2,3} {
            \draw[thick, fill=green!15] (-1.2,\i*0.8) rectangle (1.2,\i*0.8+0.7);
        }
        \node[circle, fill=red, inner sep=2pt, label=right:{$f_0(x)$}] (x2) at (0,1.4) {};
    \end{scope}

    % 可換図式
    \draw[->, thick, dashed] (x1) -- (x2);

    \node[font=\small, align=center] at (0,-0.5) {
        層構造を保つ写像
    };
\end{tikzpicture}
\end{center}

\section{まとめと発展}

\subsection{本稿で学んだこと}

\begin{enumerate}
    \item \textbf{構造塔の普遍性}
    \begin{itemize}
        \item 位相、代数、順序など多様な数学対象を統一的に扱える
        \item Bourbakiの母構造理論の現代的形式化
    \end{itemize}

    \item \textbf{豊富な具体例}
    \begin{itemize}
        \item 離散位相、位相の細かさ
        \item 主イデアル、部分群
        \item 下集合、フィルター
        \item 順序位相空間、位相群
    \end{itemize}

    \item \textbf{形式化の重要性}
    \begin{itemize}
        \item Lean 4による機械的検証可能な定義
        \item 抽象概念と具体例の橋渡し
    \end{itemize}
\end{enumerate}

\subsection{さらなる学習のために}

\begin{enumerate}
    \item \textbf{構造塔の圏論}
    \begin{itemize}
        \item 構造塔とその射は圏を成す
        \item 関手、自然変換との関係
        \item 極限・余極限の構成
    \end{itemize}

    \item \textbf{高度な具体例}
    \begin{itemize}
        \item 測度論：$\sigma$-代数の階層
        \item 確率論：フィルトレーション
        \item 代数幾何：層の構造塔
        \item ホモトピー論：CW複体の骨格
    \end{itemize}

    \item \textbf{応用}
    \begin{itemize}
        \item プログラミング言語の型システム
        \item データベースのスキーマ階層
        \item 機械学習の特徴階層
    \end{itemize}
\end{enumerate}

\begin{keyidea}
構造塔は単なる抽象理論ではなく、数学全体を統一する「言語」である。具体例を通じて理解することで、この強力な枠組みを自在に使えるようになる。
\end{keyidea}

\section*{謝辞}

本稿の形式化は OpenAI の CODEX によって、\LaTeX による解説は Anthropic の Claude Code によって生成されました。構造塔の理論は Nicolas Bourbaki の構造主義的数学観に深く影響を受けています。

\begin{thebibliography}{9}
\bibitem{bourbaki}
N. Bourbaki, \textit{Theory of Sets}, Springer, 2004.

\bibitem{bourbaki-structures}
N. Bourbaki, \textit{Elements of Mathematics: Theory of Structures},
Hermann, 1957.

\bibitem{lean}
The Lean Community, \textit{Theorem Proving in Lean 4},
\url{https://leanprover.github.io/theorem_proving_in_lean4/}

\bibitem{mathlib}
The mathlib Community, \textit{The Lean Mathematical Library},
\url{https://github.com/leanprover-community/mathlib4}

\bibitem{maclane}
S. Mac Lane, \textit{Categories for the Working Mathematician},
Springer, 1971.

\bibitem{topology}
J. L. Kelley, \textit{General Topology},
Springer, 1975.

\bibitem{algebra}
S. Lang, \textit{Algebra},
Springer, 2002.
\end{thebibliography}

\appendix
\section{Leanコード抜粋}

\subsection{構造塔の定義}

\begin{lstlisting}[style=leanstyle, caption=構造塔の基本定義]
/-- 最小層を持つ構造塔 -/
structure StructureTowerMin (X : Type*) (I : Type*)
    [Preorder I] where
  layer : I → Set X
  covering : ∀ x : X, ∃ i : I, x ∈ layer i
  monotone : ∀ {i j : I}, i ≤ j → layer i ⊆ layer j
  minLayer : X → I
  minLayer_mem : ∀ x, x ∈ layer (minLayer x)
  minLayer_minimal : ∀ x i, x ∈ layer i → minLayer x ≤ i
\end{lstlisting}

\subsection{離散位相塔の実装}

\begin{lstlisting}[style=leanstyle, caption=離散位相塔]
def DiscreteTopologyTower (X : Type*)
    [TopologicalSpace X] :
    StructureTowerMin X (Set X) where
  layer F := F
  covering := by
    intro x
    refine ⟨{x}, by simp⟩
  monotone := by
    intro F G hFG x hx
    exact hFG hx
  minLayer := fun x => {x}
  minLayer_mem := by intro x; simp
  minLayer_minimal := by
    intro x F hx
    change {x} ⊆ F
    intro y hy
    have hy' : y = x := by simpa using hy
    simpa [hy'] using hx
\end{lstlisting}

\subsection{主イデアル塔の実装}

\begin{lstlisting}[style=leanstyle, caption=主イデアル塔]
def PrincipalIdealTower (R : Type*) [CommRing R] :
    StructureTowerMin R (Ideal R) where
  layer I := (I : Set R)
  covering := by
    intro x
    refine ⟨Ideal.span ({x} : Set R), ?_⟩
    exact Ideal.subset_span (by simp : x ∈ ({x} : Set R))
  monotone := by
    intro I J hIJ x hx
    exact hIJ hx
  minLayer := fun x => Ideal.span ({x} : Set R)
  minLayer_mem := by
    intro x
    exact Ideal.subset_span (by simp : x ∈ ({x} : Set R))
  minLayer_minimal := by
    intro x I hx
    refine Ideal.span_le.2 ?_
    intro y hy
    have hy' : y = x := by simpa using hy
    simpa [hy'] using hx
\end{lstlisting}

\end{document}
